\section{Best Time to Buy and Sell Stock II}
\label{sec:dp-stock-2}

\problemstatement{Given $\textit{price}$, find the maximum profit from
\textbf{any number} of transactions (buy then sell, buy then sell,
\ldots), as long as you are never holding more than one share at a
time (you must sell before buying again).}

\begin{patternbox}
\emph{``Unlimited transactions''} is Section~\ref{sec:dp-stocks-intro}'s
base recurrence, completely \textbf{unmodified}: selling returns to
$f(i{+}1, 0)$, leaving the door open for another buy later, exactly as
the general template already specifies.
\end{patternbox}

\subsection*{The recurrence}

\[
  f(i, \textit{holding}) =
  \begin{cases}
    \max\big(-\textit{price}[i] + f(i{+}1, 1),\ f(i{+}1, 0)\big) &
    \text{not holding}, \\
    \max\big(\textit{price}[i] + f(i{+}1, 0),\ f(i{+}1, 1)\big) &
    \text{holding.}
  \end{cases}
\]

\subsection*{Approach 1: Recursion}

\begin{lstlisting}
#include <bits/stdc++.h>
using namespace std;

int stock2Recursive(int i, int holding, vector<int>& price) {
    int n = price.size();
    if (i == n) return 0;

    if (holding) {
        return max(price[i] + stock2Recursive(i + 1, 0, price),
                   stock2Recursive(i + 1, 1, price));
    }
    return max(-price[i] + stock2Recursive(i + 1, 1, price),
               stock2Recursive(i + 1, 0, price));
}

int main() {
    vector<int> price = {7, 1, 5, 3, 6, 4};
    cout << stock2Recursive(0, 0, price) << endl; // 7
    return 0;
}
\end{lstlisting}

\begin{complexitybox}[Approach 1 Complexity]
\textbf{Time.} $O(2^n)$.
\textbf{Space.} $O(n)$ recursion stack.
\end{complexitybox}

\subsection*{Approach 2: Memoization}

\begin{lstlisting}
#include <bits/stdc++.h>
using namespace std;

int stock2Memo(int i, int holding, vector<int>& price, vector<vector<int>>& memo) {
    int n = price.size();
    if (i == n) return 0;
    if (memo[i][holding] != -1) return memo[i][holding];

    int result;
    if (holding) {
        result = max(price[i] + stock2Memo(i + 1, 0, price, memo),
                      stock2Memo(i + 1, 1, price, memo));
    } else {
        result = max(-price[i] + stock2Memo(i + 1, 1, price, memo),
                      stock2Memo(i + 1, 0, price, memo));
    }
    return memo[i][holding] = result;
}

int main() {
    vector<int> price = {7, 1, 5, 3, 6, 4};
    int n = price.size();
    vector<vector<int>> memo(n, vector<int>(2, -1));
    cout << stock2Memo(0, 0, price, memo) << endl; // 7
    return 0;
}
\end{lstlisting}

\begin{complexitybox}[Approach 2 Complexity]
\textbf{Time.} $O(n)$.
\textbf{Space.} $O(n)$ memo $+$ $O(n)$ recursion stack.
\end{complexitybox}

\subsection*{Approach 3: Tabulation}

\begin{lstlisting}
#include <bits/stdc++.h>
using namespace std;

int stock2Tabulation(vector<int>& price) {
    int n = price.size();
    vector<vector<int>> dp(n + 1, vector<int>(2, 0));

    for (int i = n - 1; i >= 0; i--) {
        dp[i][1] = max(price[i] + dp[i + 1][0], dp[i + 1][1]);
        dp[i][0] = max(-price[i] + dp[i + 1][1], dp[i + 1][0]);
    }
    return dp[0][0];
}

int main() {
    vector<int> price = {7, 1, 5, 3, 6, 4};
    cout << stock2Tabulation(price) << endl; // 7
    return 0;
}
\end{lstlisting}

\subsection*{Dry run}

\dryrun{Take $\textit{price} = [7,1,5,3,6,4]$.}

The optimal strategy buys at price $1$ (day $1$), sells at price $5$
(day $2$, profit $4$), buys again at price $3$ (day $3$), sells at
price $6$ (day $4$, profit $3$): total $4+3=7$. The tabulation table,
filled backward, accumulates exactly this multi-transaction total at
$\textit{dp}[0][0]=7$.

\begin{complexitybox}[Approach 3 Complexity]
\textbf{Time.} $O(n)$.
\textbf{Space.} $O(n)$ table, no recursion stack.
\end{complexitybox}

\subsection*{Approach 4: Space Optimization}

\begin{lstlisting}
#include <bits/stdc++.h>
using namespace std;

int stock2SpaceOptimized(vector<int>& price) {
    int n = price.size();
    int nextHold = 0, nextNotHold = 0;

    for (int i = n - 1; i >= 0; i--) {
        int curHold = max(price[i] + nextNotHold, nextHold);
        int curNotHold = max(-price[i] + nextHold, nextNotHold);
        nextHold = curHold;
        nextNotHold = curNotHold;
    }
    return nextNotHold;
}

int main() {
    vector<int> price = {7, 1, 5, 3, 6, 4};
    cout << stock2SpaceOptimized(price) << endl; // 7
    return 0;
}
\end{lstlisting}

\begin{complexitybox}[Approach 4 Complexity]
\textbf{Time.} $O(n)$.
\textbf{Space.} $O(1)$.
\end{complexitybox}

\begin{ideabox}[The Greedy Alternative, and Why It Works Here Specifically]
Unlimited, uncapped transactions also admit a one-line greedy solution:
sum every positive day-to-day price increase,
$\sum_i \max(0, \textit{price}[i]-\textit{price}[i-1])$. This works
because, with no cap on transaction count, capturing every single
upward price movement (buying right before each rise, selling right
after) is always achievable and always optimal -- a luxury the capped
variants (III, IV) and the constrained variants (cooldown, fee) do not
have, which is exactly why they need the full DP treatment.
\end{ideabox}

\begin{takeawaybox}
Stock II's recurrence \emph{is} Section~\ref{sec:dp-stocks-intro}'s base
template with nothing changed -- it exists in this chapter specifically
to demonstrate the unmodified template before every later section
perturbs it.
\end{takeawaybox}
